% ============================================================
% paper/sections/03_formal_definition_k_ea.tex
% Formal Definition of the Enterprise Continuum K_EA (finalized)
% + FWHQ-01 foundational hardening additions
% ============================================================

\section{Formal Definition of the Enterprise Continuum}

\subsection{Definition of $K_{\mathrm{EA}}$}

The Enterprise Continuum $K_{\mathrm{EA}}$ is defined as an
Ontology-of-Continua (OC) continuum instantiated strictly according to
the immutable Executable Theory Specification
\texttt{ETS.K\_EA.v1.1}.

Formally, the continuum is given by the tuple
\[
K_{\mathrm{EA}} :=
\left(
\Omega_{\mathrm{EA}},
\partial\Omega_{\mathrm{EA}},
A_{\mathrm{EA}},
P_{\mathrm{EA}},
\Theta_{\mathrm{EA}},
J_{\mathrm{EA}},
C_{\mathrm{EA}},
k_{\mathrm{EA}},
\tau_{\mathrm{EA}}
\right).
\]

Each component is defined exclusively by projection from the ETS.
No additional primitives, relations, operators, or interpretive
structure are introduced at the level of this whitepaper.

The tuple notation is declarative rather than constructive.
It enumerates the components required for definability of the continuum
but does not assign numerical values, functional forms, or empirical
interpretations.
All formal content resides in the ETS.

\subsection{State Space and Boundary}

The admissible state space $\Omega_{\mathrm{EA}}$ is the set of all
enterprise states that satisfy the admissibility constraints encoded in
$\Theta_{\mathrm{EA}}$.

The boundary $\partial\Omega_{\mathrm{EA}}$ separates admissible from
non-admissible states.
Crossing this boundary results either in a phase transition within the
continuum or in termination of the continuum itself.

Non-emptiness of $\Omega_{\mathrm{EA}}$ is a necessary condition for the
existence of $K_{\mathrm{EA}}$ as a diagnosable object.
This condition is structural and does not depend on performance,
outcomes, or desirability of states.

\subsection{Axes and Degrees of Freedom}

The set of axes $A_{\mathrm{EA}}$ defines the admissible degrees of
freedom along which enterprise states may vary.

Axes parameterize variation within $\Omega_{\mathrm{EA}}$.
They are descriptive rather than prescriptive and encode neither target
states nor optimization criteria.
Their sole function is to render admissible variation explicit and
structurally comparable.

Axes do not define directions of improvement.
They define coordinates of admissibility.

\subsection{Potentials and Flows}

Potentials $P_{\mathrm{EA}}$ represent stored structural capacities of
the enterprise continuum.
Flows $J_{\mathrm{EA}}$ describe state transitions induced by internal
dynamics and externally imposed constraints.

No assumptions of equilibrium, efficiency, or optimality are made.
Potentials and flows are evaluated only with respect to their impact on
admissibility relative to $\Theta_{\mathrm{EA}}$.

Neither potentials nor flows are interpreted as resources, KPIs, or
control variables within this document.

\subsection{Thresholds and Admissibility}

The threshold set $\Theta_{\mathrm{EA}}$ encodes admissibility
conditions through component functions $f_k$ as specified in
\texttt{ETS.K\_EA.v1.1}.

Thresholds define limits beyond which:
\begin{itemize}
  \item the continuum undergoes a phase transition,
  \item structural persistence is lost,
  \item or diagnosis becomes undefined.
\end{itemize}

Thresholds are expressed as inequalities.
No numerical values are fixed at the theory level; any values supplied
externally constitute instantiations rather than definitions.

Thresholds delimit existence and diagnosability.
They do not rank states or encode preference.

\subsection{Cycles and Persistence}

The set of cycles $C_{\mathrm{EA}}$ specifies the minimal recurrent
processes required for persistence of the continuum over its internal
time scale $\tau_{\mathrm{EA}}$.

Cycles are structural requirements.
Their absence or breakdown indicates loss of persistence even when
instantaneous admissibility conditions remain momentarily satisfied.

Cycles are not operational processes or management constructs.
They denote the minimal closure conditions required for definability
over time.

\subsection{Continuity Measure}

The continuity measure $k_{\mathrm{EA}} \in [0,1]$ quantifies the degree
of structural connectivity among admissible states under the defined
dynamics.

A value of $k_{\mathrm{EA}} = 0$ constitutes a hard termination
condition.
Intermediate values indicate partial connectivity without implying
gradual degradation, recoverability, or improvement potential.

Continuity expresses the existence of admissible transitions, not their
quality.

\subsection{Internal Time Scale}

The parameter $\tau_{\mathrm{EA}}$ denotes the characteristic internal
time scale over which required cycles must close for the continuum to be
considered persistent.

No assumption is made regarding correspondence between
$\tau_{\mathrm{EA}}$ and calendar or operational time.
It is a purely internal structural parameter.

\subsection{Executable Closure}

All components of $K_{\mathrm{EA}}$ are fully specified within
\texttt{ETS.K\_EA.v1.1}.
The definition is therefore executable-closed: no external assumptions,
hidden parameters, or interpretive degrees of freedom are required to
define the continuum.

Executable closure denotes formal completeness and inspectability.
It does not imply empirical validation, scientific authority, or
predictive power.

Executable closure is a necessary condition for reproducibility,
formal inspection, and reviewer-grade validation of the specification
as a closed formal object.

%%%End of Document%%%
